\subsection{RATIONALITY FOR A LINEAR MAPPING}

DEFINITION 3. Let \(\mathbf{V}_1, \mathbf{V}_2\) be two right vector spaces over \(\mathbf{K}\) with \(\mathbf{K}'\) -structures \(\mathbf{V}_1'\) , \(\mathbf{V}_2'\) respectively. A K-linear mapping \(f: \mathbf{V}_1 \to \mathbf{V}_2\) is said to be rational over \(\mathbf{K}'\) if \(f(\mathbf{V}_1') \subset \mathbf{V}_2'\) .

If \(V_{3}\) is a third right vector space over K, with a \(K^{\prime}\) -structure \(V_{3}^{\prime}\) and a K-linear mapping \(g:V_{2}\rightarrow V_{3}\) is rational over \(K^{\prime}\) , clearly \(g\circ f:V_{1}\rightarrow V_{3}\) is rational over \(K^{\prime}\) .

PROPOSITION 3. Let \(\mathbf{V}_1, \mathbf{V}_2\) be two right vector spaces over \(\mathbf{K}\) and \(\mathbf{V}_1', \mathbf{V}_2' \mathbf{K}'\) -structures on \(\mathbf{V}_1, \mathbf{V}_2\) respectively. \(\mathbf{V}_1\) (resp. \(\mathbf{V}_2\) ) is canonically identified with \(\mathbf{V}_1' \otimes_{\mathbb{K}'} \mathbf{K}\) (resp. \(\mathbf{V}_2' \otimes_{\mathbb{K}'} \mathbf{K}\) ) (no. 1, Proposition 1).

\begin{enumerate}
\def\labelenumi{(\roman{enumi})}
\item
  The mapping \(f' \mapsto f' \otimes 1_{\mathbf{K}} = f_{(\mathbf{K})}'\) is a bijection of \(\operatorname{Hom}_{\mathbf{K}'}(\mathbf{V}_1', \mathbf{V}_2')\) onto the set of K-linear mappings of \(\mathbf{V}_1\) into \(\mathbf{V}_2\) which are rational over \(\mathbf{K}'\) ; the inverse bijection associates with every K-linear mapping \(f: \mathbf{V}_1 \to \mathbf{V}_2\) rational over \(\mathbf{K}'\) the K'-linear mapping \(f': \mathbf{V}_1' \to \mathbf{V}_2'\) with the same graph as the restriction off to \(\mathbf{V}_1'\) .
\item
  For every K-linear mapping \(f: \mathbf{V}_1 \to \mathbf{V}_2\) rational over \(\mathbf{K}'\) ,
\end{enumerate}

\[
f (\mathrm{V} _ {1} ^ {\prime}) = f (\mathrm{V} _ {1}) \cap \mathrm{V} _ {2} ^ {\prime} \quad and \quad \bar {f} ^ {- 1} (\mathrm{V} _ {2} ^ {\prime}) = \mathrm{V} _ {1} ^ {\prime} + \operatorname{Ker} (f).
\]

\begin{enumerate}
\def\labelenumi{(\roman{enumi})}
\item
  Clearly with the above identifications, if \(f': \mathbf{V}_1' \to \mathbf{V}_2'\) is a K'-linear mapping, \(f_{(\mathbf{K})}' = f' \otimes 1_{\mathbf{K}}\) is rational over \(\mathbf{K}'\) and \(f'\) is the mapping with the same graph as the restriction of \(f_{(\mathbf{K})}'\) to \(\mathbf{V}_1'\) . Conversely, if \(f: \mathbf{V}_1 \to \mathbf{V}_2\) is a K-linear mapping rational over \(\mathbf{K}'\) and \(f': \mathbf{V}_1' \to \mathbf{V}_2'\) has the same graph as the restriction of \(f\) to \(\mathbf{V}_1'\) , \(f\) and \(f_{(\mathbf{K})}'\) coincide on \(\mathbf{V}_1'\) , which is a generating system of \(\mathbf{V}_1\) over \(\mathbf{K}\) , hence \(f = f_{(\mathbf{K})}'\) .
\item
  If \(f = f' \otimes 1_{\mathbb{K}}\) , then \(f(\mathbf{V}_1) = f(\mathbf{V}_1' \otimes_{\mathbb{K}'} \mathbf{K}) = f'(\mathbf{V}_1') \otimes_{\mathbb{K}'} \mathbf{K}\) and, as \(f'(\mathbf{V}_1') \subset \mathbf{V}_2'\) , \(f(\mathbf{V}_1') = f'(\mathbf{V}_1') = f(\mathbf{V}_1) \cap \mathbf{V}_2'\) (§ 7, no. 9, Proposition 19); the formula \(f^{-1}(\mathbf{V}_2') = \mathbf{V}_1' + \operatorname{Ker}(f)\) follows immediately.
\end{enumerate}

COROLLARY 1. In the notation of Proposition 3, \(\operatorname{Im}(f) = (\operatorname{Im}(f'))_{(\mathbb{K})}\) ,

\[
\operatorname{Ker} (f) = (\operatorname{Ker} (f ^ {\prime})) _ {(\mathbb {K})}, \quad \operatorname{Coker} (f) = (\operatorname{Coker} (f ^ {\prime})) _ {(\mathbb {K})}.
\]

In particular, for \(f\) to be injective (resp. surjective, zero), it is necessary and sufficient that \(f'\) be so. If \(f\) is bijective, its inverse mapping is rational over \(\mathbf{K}'\) .

This is a particular case of § 7, no. 9, Corollary to Proposition 18.

COROLLARY 2. Let \(f: \mathbf{V}_1 \to \mathbf{V}_2\) be a K-linear mapping rational over \(\mathbf{K}'\) . For every vector sub-K-space \(\mathbf{W}_1\) of \(\mathbf{V}_1\) (resp. \(\mathbf{W}_2\) of \(\mathbf{V}_2\) ) rational over \(\mathbf{K}'\) , \(f(\mathbf{W}_1)\) (resp. \(f(\mathbf{W}_2)\) ) is a vector sub-K-space of \(\mathbf{V}_2\) (resp. \(\mathbf{V}_1\) ) rational over \(\mathbf{K}'\) .

In the notation of Proposition 3, for every vector sub-K'-space W \(_{1}\) of V \(_{1}\) , f \(_{(K)}\) (W \(_{1}\) ' ⊗ \(_{K'}\) K) = f'(W \(_{1}\) ) ⊗ \(_{K'}\) K; whence the assertion relating to W \(_{1}\) (no. 2, Proposition 2). On the other hand, let W \(_{2}\) be a vector sub-K'-space of V \(_{2}\) and let g' be the canonical K'-linear mapping V \(_{2}\) → V \(_{2}\) /W \(_{2}\) ; then \(f'^{-1}(W_{2}') = \text{Ker}(g' \circ f')\) ; hence, by Corollary 1,

\[
\bar {f} _ {(\mathbf {K})} ^ {\prime - 1} (\mathbf {W} _ {2} ^ {\prime} \otimes_ {\mathbf {K} ^ {\prime}} \mathbf {K}) = \bar {f} ^ {\prime - 1} (\mathbf {W} _ {2} ^ {\prime}) \otimes_ {\mathbf {K} ^ {\prime}} \mathbf {K},
\]

whence the assertion relating to \(W_{2}\) .

Let \(V_{1}, V_{2}\) be two right vector K-spaces with \(K'\) -structures \(V_{1}', V_{2}'\) respectively. It is immediate that \(V_{1}' \times V_{2}'\) is a \(K'\) -structure on \(V_{1} \times V_{2}\) , called the product of the \(K'\) -structures \(V_{1}'\) and \(V_{2}'\) .

PROPOSITION 4. For a K-linear mapping \(f: \mathbf{V}_1 \to \mathbf{V}_2\) to be rational over \(\mathbf{K}'\) , it is necessary and sufficient that its graph \(\Gamma\) be rational over \(\mathbf{K}'\) for the product \(\mathbf{K}'\) -structure on \(\mathbf{V}_1 \times \mathbf{V}_2\) .

Let \(g\) be the mapping \(x_1 \mapsto (x_1, f(x_1))\) from \(V_1\) to \(V_1 \times V_2\) ; this is a K-linear mapping such that \(\Gamma = g(V_1)\) ; if \(f\) is rational over \(K'\) , so is \(\Gamma\) by virtue of Corollary 2 to Proposition 3. Conversely, suppose that \(\Gamma\) is rational over \(K'\) and let it have the \(K'\) -structure induced by that on \(V_1 \times V_2\) ; it follows immediately from the definitions that the restrictions \(p_1, p_2\) to \(\Gamma\) of the projections \(pr_1, pr_2\) , are K-linear mappings, rational over \(K'\) of \(\Gamma\) into \(V_1\) and \(V_2\) respectively. As \(p_1\) is bijective, its inverse mapping \(q_1\) is rational over \(K'\) (Corollary 1 to Proposition 3) and hence so is \(f = p_2 \circ q_1\) .

